\section*{NeurIPS Paper Checklist}

\begin{enumerate}

\item {\bf Claims}
    \item[] Question: Do the main claims made in the abstract and introduction accurately reflect the paper's contributions and scope?
    \item[] Answer: \answerYes{}
    \item[] Justification: The five contributions bulleted in \S\ref{sec:intro} (cascade router, personalization-by-compression, bounded conversation buffer, routing log with CSV export, open-source implementation with a passing test suite) each have a corresponding section in the body (\S\ref{sec:router}, \S\ref{sec:compress}, \S\ref{sec:buffer}, \S\ref{sec:log}, \S\ref{sec:experiments}--\S\ref{sec:results}). We do not claim performance beyond the 22-example held-out split (Table~\ref{tab:router}).

\item {\bf Limitations}
    \item[] Question: Does the paper discuss the limitations of the work performed by the authors?
    \item[] Answer: \answerYes{}
    \item[] Justification: \S\ref{sec:limits} discusses seed-set scale and the weak Wilson lower bound at 22 examples, the binary-only label space, summarization fidelity risks, residual profile leakage to the small LLM, and the fact that $\tau$ was chosen by eye rather than calibrated.

\item {\bf Theory assumptions and proofs}
    \item[] Question: For each theoretical result, does the paper provide the full set of assumptions and a complete (and correct) proof?
    \item[] Answer: \answerNA{}
    \item[] Justification: The paper presents a system design and an empirical evaluation; it contains no theorems requiring proof. Equations~\ref{eq:cascade} and \ref{eq:summary} are definitions of the dispatch rule and incremental update, not claims.

\item {\bf Experimental result reproducibility}
    \item[] Question: Does the paper fully disclose all the information needed to reproduce the main experimental results of the paper to the extent that it affects the main claims and/or conclusions of the paper (regardless of whether the code and data are provided or not)?
    \item[] Answer: \answerYes{}
    \item[] Justification: Appendix~\ref{app:repro} gives the exact commands (\texttt{python -m router.train}, \texttt{pytest -q}); the seed data lives in \texttt{router/seed\_data.py} inside the released repo, and all hyperparameters (TF-IDF \texttt{ngram\_range}, logistic regression \texttt{C}, split seed, thresholds) are listed in Appendix~B.

\item {\bf Open access to data and code}
    \item[] Question: Does the paper provide open access to the data and code, with sufficient instructions to faithfully reproduce the main experimental results, as described in supplemental material?
    \item[] Answer: \answerYes{}
    \item[] Justification: Code, seed data, trained v0 model weights, and tests are publicly available at \url{https://github.com/sandeepvijayarao09/dual-llm-system} under an open-source license; Appendix~\ref{app:repro} gives a three-command reproduction path.

\item {\bf Experimental setting/details}
    \item[] Question: Does the paper specify all the training and test details (e.g., data splits, hyperparameters, how they were chosen, type of optimizer) necessary to understand the results?
    \item[] Answer: \answerYes{}
    \item[] Justification: \S\ref{sec:router} specifies the feature set and the full sklearn pipeline; \S\ref{sec:experiments} specifies the 80/20 stratified split; Appendix~B lists every threshold and hyperparameter used.

\item {\bf Experiment statistical significance}
    \item[] Question: Does the paper report error bars suitably and correctly defined or other appropriate information about the statistical significance of the experiments?
    \item[] Answer: \answerNo{}
    \item[] Justification: We report a single stratified split result rather than cross-validated error bars because the held-out set has only 22 examples; in \S\ref{sec:limits} we explicitly note the Wilson 95\% lower bound and flag this as the motivation for the routing-log-driven retraining loop.

\item {\bf Experiments compute resources}
    \item[] Question: For each experiment, does the paper provide sufficient information on the computer resources (type of compute workers, memory, time of execution) needed to reproduce the experiments?
    \item[] Answer: \answerYes{}
    \item[] Justification: \S\ref{sec:experiments} states that all experiments run on a single MacBook Pro (Apple M-series CPU, 16\,GB RAM), that \texttt{pytest -q} completes in under 5 seconds, and that router training completes in under 1 second.

\item {\bf Code of ethics}
    \item[] Question: Does the research conducted in the paper conform, in every respect, with the NeurIPS Code of Ethics \url{https://neurips.cc/public/EthicsGuidelines}?
    \item[] Answer: \answerYes{}
    \item[] Justification: The project uses only synthetic/hand-authored queries, no human-subject data, and public commercial APIs under their stated terms of service; the authors reviewed the NeurIPS Code of Ethics and found no conflicts.

\item {\bf Broader impacts}
    \item[] Question: Does the paper discuss both potential positive societal impacts and negative societal impacts of the work performed?
    \item[] Answer: \answerYes{}
    \item[] Justification: \S\ref{sec:impacts} discusses energy/cost reductions on the positive side and silent quality downgrades plus routing-log PII risk on the negative side.

\item {\bf Safeguards}
    \item[] Question: Does the paper describe safeguards that have been put in place for responsible release of data or models that have a high risk for misuse (e.g., pre-trained language models, image generators, or scraped datasets)?
    \item[] Answer: \answerNA{}
    \item[] Justification: The released router is a TF-IDF + logistic regression classifier trained on 108 benign hand-authored queries; it cannot generate text and poses no dual-use risk. The underlying LLMs (GPT-4o, GPT-4o-mini) are accessed via OpenAI's API under OpenAI's existing safeguards and are not released by us.

\item {\bf Licenses for existing assets}
    \item[] Question: Are the creators or original owners of assets (e.g., code, data, models), used in the paper, properly credited and are the license and terms of use explicitly mentioned and properly respected?
    \item[] Answer: \answerYes{}
    \item[] Justification: scikit-learn (BSD-3), joblib (BSD-3), pytest (MIT), Streamlit (Apache 2.0), python-dotenv (BSD-3), and the OpenAI Python SDK (Apache 2.0) are all used under their original licenses and cited in the references or the project \texttt{requirements.txt}. GPT-4o and GPT-4o-mini are used via the OpenAI API under OpenAI's API Terms of Use \citep{openai_pricing_2024}.

\item {\bf New assets}
    \item[] Question: Are new assets introduced in the paper well documented and is the documentation provided alongside the assets?
    \item[] Answer: \answerYes{}
    \item[] Justification: The released repository includes a \texttt{README.md}, per-module docstrings, typed dataclasses for the router's and orchestrator's return values, a trained \texttt{router\_v0.joblib} weights file, and a 38-test pytest suite that doubles as executable documentation.

\item {\bf Crowdsourcing and research with human subjects}
    \item[] Question: For crowdsourcing experiments and research with human subjects, does the paper include the full text of instructions given to participants and screenshots, if applicable, as well as details about compensation (if any)?
    \item[] Answer: \answerNA{}
    \item[] Justification: The paper involves no crowdsourcing and no human-subjects research; the 108-example seed was authored by the three authors themselves.

\item {\bf Institutional review board (IRB) approvals or equivalent for research with human subjects}
    \item[] Question: Does the paper describe potential risks incurred by study participants, whether such risks were disclosed to the subjects, and whether Institutional Review Board (IRB) approvals (or an equivalent approval/review based on the requirements of your country or institution) were obtained?
    \item[] Answer: \answerNA{}
    \item[] Justification: No human subjects, crowdworkers, or external data collection were involved, so IRB review does not apply.

\item {\bf Declaration of LLM usage}
    \item[] Question: Does the paper describe the usage of LLMs if it is an important, original, or non-standard component of the core methods in this research? Note that if the LLM is used only for writing, editing, or formatting purposes and does \emph{not} impact the core methodology, scientific rigor, or originality of the research, declaration is not required.
    \item[] Answer: \answerYes{}
    \item[] Justification: LLMs are a core component of the system under study: GPT-4o-mini is used as the small-route answerer, LLM classifier, summarizer, personalizer, and conversation-buffer compressor; GPT-4o is used as the big-route reasoner. Their roles are described in detail in \S\ref{sec:system}--\S\ref{sec:method} and summarized in Table~\ref{tab:components}. The authors additionally used an LLM-based coding assistant for code drafting and manuscript copyediting; it was not used for experimental design, analysis, or result generation.

\end{enumerate}
